\chapter{Auditing Compatibility with xna4-spec}
\label{ch:xna4-spec-auditing}

Every coverage percentage this book has cited --- 227 of 245 public FNA types, 92.7\%,
Chapter~\ref{ch:what-is-cna} --- came from comparing CNA against a ground truth rather than
against memory or against FNA's own source. This chapter covers that ground truth directly:
\texttt{xna4-spec}, and what auditing against it actually looks like.

\section{What xna4-spec actually is}

\texttt{xna4-spec} is a complete, machine-readable XML transcription of the official XNA~4.0
documentation, originally published at Microsoft's own
\texttt{learn.microsoft.com/\allowbreak{}.../\allowbreak{}previous-versions/\allowbreak{}windows/\allowbreak{}xna/}
archive --- every class,
struct, enum, interface, and delegate across all nineteen XNA namespaces, 544 types in total,
with each type's original property, method, constructor, field, and event descriptions
preserved. A master XSD schema and a top-level index of all 544 types organize the whole
collection; one directory exists per namespace, mirroring the real
\texttt{Microsoft.Xna.Framework.*} namespace tree exactly.

Per-namespace type counts give a sense of where the real API surface concentrates:
\cnans{Microsoft.Xna.Framework.Graphics} alone accounts for 175 of the 544 types, with
\cnans{GamerServices} (53) and \cnans{Content.Pipeline.Graphics} (46) the next largest ---
which is exactly why this book devoted an entire Part to Graphics and this Part gives
GamerServices its own chapter, rather than the namespace breakdown being an arbitrary
organizing choice on this book's part.

\section{Why this matters more than diffing against FNA}

Chapter~\ref{ch:media-system} already gave the sharpest illustration of why this distinction
is not academic: eight separate audit rounds of CNA's \cnaclass{Song} class diffed against
FNA's own source and missed a real gap (FNA itself omits \texttt{Album}, \texttt{Artist},
\texttt{Genre}, and \texttt{ToString()} relative to the real reference assemblies), because
every round trusted FNA as a stand-in for the real specification rather than checking the
real specification directly. \texttt{xna4-spec} exists precisely to make that class of
mistake avoidable: a diff against it checks CNA against the actual, original Microsoft
documentation, with FNA's own gaps and quirks no longer able to hide a CNA gap behind them.

\section{Running an audit yourself}

The practical shape of an audit, drawn from how CNA's own coverage figures are produced
(Chapter~\ref{ch:project-practice}'s \texttt{AUDIT.md} conventions), is a straightforward
three-step comparison: pick a namespace or class from \texttt{xna4-spec}'s own index, read
its full member list (properties, methods, constructors, fields, events) from the
corresponding XML file, and check each member off against CNA's actual header for that class
--- present and correctly shaped, present but tagged \texttt{NOXNA} (and therefore not really
part of what you are checking), or missing. A member present in CNA but absent from
\texttt{xna4-spec} and not tagged \texttt{NOXNA} at all is the specific ``Extra-unmarked''
finding category Chapter~\ref{ch:devices-sensors} showed a real instance of
(\cnaclass{SensorBase}'s \texttt{TimeBetweenUpdatesChanged} event, found and then correctly
retagged) --- exactly the class of bug the \texttt{NOXNA} compile-time gate from
Chapter~\ref{ch:testing-philosophy} exists to catch mechanically, with a manual
\texttt{xna4-spec} audit as the complementary, judgment-requiring check for anything that
gate cannot see on its own (a member that is tagged correctly but has the wrong signature, for
instance).

\section{A worked audit: \cnaclass{Ray}, member by member}

The three-step process above is easiest to see against a small, complete type rather than a
large one. \texttt{xna4-spec}'s own
\texttt{Microsoft.Xna.Framework/Ray.xml} lists, in full: two fields (\texttt{Position},
\texttt{Direction}, both \texttt{Vector3}), one constructor (\texttt{Ray(Vector3, Vector3)}),
and six methods --- \texttt{Equals} (two overloads, one taking a boxed \texttt{Object} and one
taking another \texttt{Ray}), \texttt{GetHashCode}, \texttt{ToString}, the \texttt{==} / \texttt{!=}
operators, and \texttt{Intersects} against each of \cnaclass{BoundingBox}, \cnaclass{Plane},
\cnaclass{BoundingSphere}, and \cnaclass{BoundingFrustum} (the first three of those each with a
second, \texttt{ref} / \texttt{out} overload; \texttt{BoundingFrustum} has only the one
return-value form in the real spec, not a \texttt{ref} / \texttt{out} pair, which matters below).
Checking that list member-by-member against \cnaclass{Ray}'s real header
(\texttt{include/Microsoft/Xna/Framework/Ray.hpp}) finds a complete, correctly-shaped match for
every one of them: both fields present with matching types; the two-argument constructor
present; all four \texttt{Intersects} targets present, each with exactly the same
value-returning/output-parameter overload pattern the spec lists --- including the detail that
\texttt{BoundingFrustum} correctly has only the single return-value overload, with no spurious
\texttt{ref} / \texttt{out} pair invented for it; \texttt{ToString}, and the free-function
\texttt{operator==} / \texttt{operator!=} pair implementing \texttt{op\_Equality} /
\texttt{op\_Inequality} exactly as this book's naming conventions predict.

Two small, precise deviations are worth naming exactly rather than glossing over, because they
illustrate two different categories an audit needs to tell apart. First, \cnaclass{Ray} has no
\texttt{Equals(Object)} overload at all --- and this is correctly not a gap: C\texttt{++} has no
universal boxed root type the way every C\# value type inherits from \texttt{System.Object}, so
there is no C\texttt{++} member this overload could ever correspond to. This is a third audit
outcome distinct from both ``missing'' and ``present but \texttt{NOXNA}-tagged'': a member that
cannot be translated at all because the language feature it depends on does not exist on the
C\texttt{++} side, and is therefore correctly absent with no tag needed to explain the absence.
Second, \texttt{GetHashCode} returns \texttt{std::size\_t} in \cnaclass{Ray} rather than the
spec's \texttt{int} --- a deliberate, sensible convention substitution (matching how
C\texttt{++}'s own \texttt{std::hash} ecosystem represents a hash value) rather than a
signature bug, the same class of considered adaptation Chapter~\ref{ch:migration-guide}
described for the \texttt{getXProperty()} / \texttt{setXProperty()} translation of C\# properties.
\cnaclass{Ray} also adds one member \texttt{xna4-spec} does not list at all: a defaulted
zero-argument constructor (\texttt{Ray() = default;}). This is not flagged as an
``Extra-unmarked'' finding in the sense Chapter~\ref{ch:devices-sensors}'s
\cnaclass{SensorBase} case was --- default-constructibility is a baseline C\texttt{++}
expectation for a plain aggregate struct, not a deliberate extension of the XNA surface, and
this book's own audits do not require a \texttt{NOXNA} tag for it.

\section{A second worked audit: \cnaclass{Viewport}, member by member}

Every example above comes from \texttt{Microsoft.Xna.Framework}, the root namespace. The same
three-step process applies unchanged to any of the other eighteen namespaces
\texttt{xna4-spec} covers --- worth demonstrating once with a type from elsewhere in the tree,
both to confirm the process generalizes and because this second audit turns up a genuinely
different kind of finding than \cnaclass{Ray} did.

\cnaclass{Viewport} lives in \cnans{Microsoft.Xna.Framework.Graphics}, and \texttt{xna4-spec}'s
own \texttt{Microsoft.Xna.Framework.Graphics/Viewport.xml} lists: two constructors
(\texttt{Viewport(Rectangle)} and \texttt{Viewport(Int32, Int32, Int32, Int32)}), nine
properties --- \texttt{AspectRatio} (get-only), \texttt{Bounds}, \texttt{Height},
\texttt{MaxDepth}, \texttt{MinDepth}, \texttt{TitleSafeArea} (get-only), \texttt{Width},
\texttt{X}, and \texttt{Y} --- and three methods: \texttt{Project}, \texttt{Unproject}, and
\texttt{ToString}. No fields and no events are declared at all, matching \cnaclass{Viewport}'s
own nature as a pure value-type view onto four coordinates and a depth range.

Checking that list against the real header
(\texttt{include/\allowbreak{}Microsoft/\allowbreak{}Xna/\allowbreak{}Framework/\allowbreak{}Graphics/\allowbreak{}Viewport.hpp})
finds every member present. Both constructors match exactly, including the same defaulted
\texttt{MinDepth=0} / \texttt{MaxDepth=1} the spec's own remarks describe for each --- and, as
with \cnaclass{Ray}, \cnaclass{Viewport} also adds one member the spec does not list: a
defaulted zero-argument constructor, the same baseline-aggregate case already established as
not requiring a \texttt{NOXNA} tag. \texttt{Height}, \texttt{MaxDepth}, \texttt{MinDepth},
\texttt{Width}, \texttt{X}, and \texttt{Y} are each backed by \texttt{DEF\_PROP} --- the same
getter/setter-generating macro whose \texttt{getXProperty()} / \texttt{setXProperty()} naming
convention Chapter~\ref{ch:migration-guide} already covers --- while \texttt{AspectRatio},
\texttt{Bounds}, and \texttt{TitleSafeArea} are hand-written methods instead. That split is not
an inconsistency; it tracks a real distinction in what each property actually is.
\texttt{DEF\_PROP} models a property backed directly by one stored field; \texttt{AspectRatio}
is computed from two fields with no field of its own, \texttt{Bounds} constructs a
\cnaclass{Rectangle} from four fields (and, being settable, unpacks one back into all four on
the way in), and \texttt{TitleSafeArea} --- see below --- is computed from \texttt{Bounds}
again. None of the three can be expressed as a direct getter/setter pair over a single field,
so all three fall outside what \texttt{DEF\_PROP} can generate and are written out by hand,
exactly as the macro's own design would predict once you know what it does and does not cover.
\texttt{Project} and \texttt{Unproject} both match the spec's four-parameter
\texttt{(Vector3, Matrix, Matrix, Matrix)} shape exactly, confirmed directly against
\texttt{Viewport.cpp}'s own implementation of each.

The audit's one substantive finding sits inside that hand-written trio. \texttt{xna4-spec}'s
own remarks for \texttt{TitleSafeArea} are specific about what the property is supposed to do:
the title safe area is ``the part of the screen where text should be displayed,'' and the
property itself ``reduces the title safe area of the viewport uniformly.'' A real console-era
XNA implementation insets the returned rectangle by a fixed margin (historically an
80\%-of-viewport safe region on Xbox 360, the platform the remark is written for).
\cnaclass{Viewport}'s own header describes its C\texttt{++} member differently --- ``Same
rectangle as Bounds'' --- and \texttt{getTitleSafeAreaProperty()}'s real implementation in
\texttt{Viewport.cpp} bears that out exactly: \texttt{return getBoundsProperty();}, no inset
applied at all. This is not an untested oversight either; \texttt{ViewportTests.cpp} carries a
dedicated \texttt{TitleSafeAreaEqualsBounds} test that asserts precisely this behavior ---
\texttt{TitleSafeArea} returning the same four coordinates as \texttt{Bounds} is the
deliberate, regression-tested baseline, not a bug nobody has gotten to yet.

By the three-step process alone, \texttt{TitleSafeArea} is a clean pass: present, correctly
named, correctly typed (\cnaclass{Rectangle}), get-only exactly as the spec describes, no
\texttt{NOXNA} tag needed because nothing about it is an addition. It is the first concrete
instance in this chapter of the limitation the closing section below names in the abstract: a
membership-and-shape audit has nothing further to say once a member exists and type-checks,
even when --- as here --- its actual behavior openly diverges from what the specification's
own prose says that member is for. Whether the divergence is the right call is a separate,
judgment-requiring question (the physical off-screen bleed an 80\%-safe-area convention
corrected for on a CRT-era television essentially does not exist on a modern flat-panel
display, which is a real argument for leaving \texttt{Viewport} as-is) --- but that judgment is
not something the audit itself can render; it can only surface the gap between documented and
implemented behavior for a human to weigh next.

A practical heuristic falls out of how this specific finding was actually found: the properties
worth reading past their signature are the hand-written ones, not the \texttt{DEF\_PROP} ones.
A \texttt{DEF\_PROP}-backed property is a direct read or write of one stored field --- it either
holds the value it was given or it does not, and \cnaclass{Ray}'s own audit above already
establishes that class of check is straightforward. A hand-written accessor is where an
implementation choice can hide, because it is free to compute something other than what its
name and the spec's remarks describe. When auditing a type with a mix of both, reading every
hand-written property's full \texttt{<remarks>} block in \texttt{xna4-spec}'s own XML --- not
just its signature --- is the specific step that catches a case like \texttt{TitleSafeArea}.

\section{A licensing note worth stating plainly}

\texttt{xna4-spec}'s own README is direct about its content's origin: the API names,
descriptions, and everything else in the collection are derived from Microsoft's own XNA~4.0
documentation and remain Microsoft's intellectual property. The repository presents itself as
an unofficial, community-maintained conversion into a machine-readable format, intended
solely for non-commercial, informational, and educational use --- the same spirit in which
this book cites specific facts drawn from it.

\section{What an audit cannot tell you}

One limit is worth naming directly: \texttt{xna4-spec} tells you whether a member exists and
what it is documented to do --- it cannot tell you whether CNA's own implementation of a
present, correctly-shaped member actually behaves correctly. That is exactly the job the
pixel tests, differential testing, and oracle methodology from
Chapter~\ref{ch:testing-philosophy} exist to do instead. A useful mental model for the rest of
this Part: \texttt{xna4-spec} auditing answers ``is the surface complete and correctly
shaped,'' while the testing mechanisms from the previous chapter answer ``does the surface
that exists actually work.'' Both questions matter, and neither answers the other.

The \cnaclass{Viewport} / \texttt{TitleSafeArea} case above is a concrete instance of exactly
this limit, with a further wrinkle worth naming: \texttt{TitleSafeAreaEqualsBounds} is a real,
passing unit test, so this is not a gap either kind of check would surface as a failure. The
test faithfully verifies the implementation's own chosen behavior; it says nothing about
whether that behavior matches \texttt{xna4-spec}'s remarks, because nothing in this project's
test suite is written against \texttt{xna4-spec} directly. Closing a gap like this one, if it
is ever judged worth closing, starts with the manual read this chapter demonstrates --- not
with a red test, because there is not one.
